% =============================================================================
% preamble_accv.tex — ACCV 版数学宏 + theorem 环境
%
% 与 preamble.tex 的差异：
%   [C1] 删除 \eg \ie \etc \etal 定义（由 accvabbrv.sty 负责，避免 \already defined 冲突）
%        保留 \aka（accvabbrv 未定义）
%   [C3] 删除 \usepackage{cleveref}（由 accv.sty \AtEndPreamble 自动加载）
%   [C2] 删除 \usepackage[round]{natbib}（已在 main_accv.tex 加载于 accv 之前）
%   [C5] \newtheorem 在 amsthm 下定义，llncs 的 \spnewtheorem 不与其冲突
%   其余数学宏与 preamble.tex 完全一致，以保证 drafts_short/ 正常编译
% =============================================================================

% ── 已在 main_accv.tex 顶层加载，此处不重复 ──────────────────────────────────
% amsmath, amssymb, mathtools, bm, graphicx, booktabs, multirow,
% xcolor[table], xspace, hyperref, natbib — 均已在 main_accv.tex 处理
%
% NOTE: amsthm 不在此加载。llncs.cls 已通过 \spnewtheorem / \spn@wtheorem
%   定义 theorem / lemma / corollary / proposition / remark / proof / claim 等。
%   在 llncs 下再 \usepackage{amsthm} + \newtheorem 会触发
%   "Command \theorem already defined" fatal error [C5]。

% ── \aka（accvabbrv 未提供的缩写）[C1] ──────────────────────────────────────
\providecommand{\aka}{a.k.a.\xspace}

% ── Theorem 环境 [C5] ────────────────────────────────────────────────────────
% theorem / lemma / corollary / proposition / remark / proof / claim
%   已由 llncs.cls 的 \spn@wtheorem 定义，此处不重复。
%
% assumption：llncs 未定义，用 \spnewtheorem 补充
\spnewtheorem{assumption}{Assumption}{\bfseries}{\itshape}

% ── \qedhere 兼容 [C5-补] ────────────────────────────────────────────────────
% amsthm 的 \qedhere 在 llncs proof 环境内无效（llncs 不加载 amsthm）。
% 用 \providecommand 静默忽略；llncs proof 环境末尾已自动插入 QED 方框。
\providecommand{\qedhere}{}

% ── 匿名模型名宏（de-anon 只需改 RHS）──────────────────────────────────────
\newcommand{\qvib}{QC-VIB}
\newcommand{\visiscore}{5-head IQA}
\newcommand{\visienhance}{QP-Enhance}
\newcommand{\qcts}{QC-TS}
\newcommand{\itb}{QSB}

% ── 数学算子 ────────────────────────────────────────────────────────────────
\DeclareMathOperator{\KL}{KL}
\DeclareMathOperator*{\E}{\mathbb{E}}
\DeclareMathOperator{\Var}{Var}
\DeclareMathOperator{\softmax}{softmax}
\DeclareMathOperator{\sigmoid}{sigmoid}
\DeclareMathOperator{\tr}{tr}
\DeclareMathOperator*{\argmax}{arg\,max}
\newcommand{\TV}{\mathrm{TV}}
\newcommand{\Hent}{H}
\newcommand{\qbar}{\bar{q}}
\newcommand{\Tw}{T_\omega}
\newcommand{\Zenh}{Z_{\mathrm{enh}}}
\newcommand{\Zref}{Z_{\mathrm{ref}}}
\newcommand{\Lq}{L_{q_\theta}}
\newcommand{\Ldp}{\mathcal{L}_{\mathrm{DP}}}
\newcommand{\Ragent}{\mathcal{R}_{\mathrm{agent}}}
\newcommand{\Rdirect}{\mathcal{R}_{\mathrm{direct}}}
\newcommand{\Renh}{\mathcal{R}_{\mathrm{enh}}}
\newcommand{\tauenh}{\tau_{\mathrm{enh}}}
\newcommand{\tauhigh}{\tau_{\mathrm{high}}}
\newcommand{\taulow}{\tau_{\mathrm{low}}}
% Theorem 2 (agent risk bound)
\newcommand{\Gen}{G_{\mathrm{enh}}}
\newcommand{\Gret}{G_{\mathrm{retake}}}
\newcommand{\pisalv}{\pi_{\mathrm{salvage}}}
\newcommand{\hhigh}{h_{\mathrm{high}}}
\newcommand{\actd}{a_d}
\newcommand{\acte}{a_e}
\newcommand{\actq}{a_q}
\newcommand{\actr}{a_r}
% Corollary 1 (composite ECE bound)
\DeclareMathOperator{\ECE}{ECE}
\newcommand{\Lt}{L_T}
\newcommand{\Tmin}{T_{\min}}
\newcommand{\epsqts}{\epsilon_{\mathrm{qts}}}
\newcommand{\pQV}{\hat{p}^{(\mathrm{QV})}}
\newcommand{\pQCTS}{\hat{p}^{(\mathrm{QCTS})}}
\newcommand{\pComp}{\hat{p}^{(\mathrm{comp})}}

% ── Draft-stage TODO marker（红色 [TODO] 占位，不输出参数文字）────────────────
\newcommand{\todo}[1]{\textcolor{red}{\textbf{[TODO]}}}
